\hypertarget{appendix-173b52c063ae}{}
\section{Vector spaces subspaces and span}\label{appendix-173b52c063ae-vector-spaces-subspaces-and-span}

\textbf{Curriculum depth:} Introduced · \textbf{First seen:} Phase6 · \textbf{Latest development:} Phase6

\subsection{Definition and mechanism}\label{appendix-173b52c063ae-definition-and-mechanism}

A vector space is closed under addition and field-scalar multiplication with coherent algebraic laws. A subspace retains these operations inside a larger space. The span of generators is every finite linear combination of them.

The plane \(z=0\) is a subspace of \(\mathbb R^3\); \(z=1\) is an affine translate, not a subspace. The span of \((4,0,0)\) and \((0,3,0)\) contains all in-plane movements.

\subsection{Prerequisites}\label{appendix-173b52c063ae-prerequisites}

\begin{itemize}
\tightlist
\item
  \hyperref[appendix-0659864c0aa3]{Vectors coordinates and affine points} --- Point, displacement and coordinate roles determine which geometric operations are meaningful.
\end{itemize}

\subsection{Used by}\label{appendix-173b52c063ae-used-by}

\begin{itemize}
\item
  \href{../01-Foundations-I/main.pdf\#nameddest=appendix-ce0341f7ea0d}{Basis coordinates and dimension (01-Foundations-I, appendix-ce0341f7ea0d)} --- Builds on this mechanism in the connected geometry story.
\item
  \href{../04-Visual-Computing-II/main.pdf\#nameddest=appendix-b0732b22e7d6}{Additive energy and approximate structure (04-Visual-Computing-II, appendix-b0732b22e7d6)}
\item
  \href{../04-Visual-Computing-II/main.pdf\#nameddest=appendix-3c73f7862157}{Exceptional projections and dimension (04-Visual-Computing-II, appendix-3c73f7862157)} \#\# Recall and next depth
\end{itemize}

Distinguish generators from their span; prove the one-step subspace criterion. Full production proofs, target-specific integration and learner mastery require further evidence.

\subsection{Source and navigation}\label{appendix-173b52c063ae-source-and-navigation}

\href{../01-Foundations-I/main.pdf\#nameddest=01-foundations-i-ch017-the-reachable-movements-form-a-vector-space}{Teaching narrative (01-Foundations-I, 01-foundations-i-ch017-the-reachable-movements-form-a-vector-space)} · \href{../01-Foundations-I/main.pdf\#nameddest=appendix-fd090d982a15-phase-6---parts-501-600}{Coverage and source gap (01-Foundations-I, appendix-fd090d982a15-phase-6-\/-\/-parts-501-600)} · \href{../01-Foundations-I/main.pdf\#nameddest=appendix-a103e7846d21-phase-6---source-analysis-and-verification}{Evidence and limits (01-Foundations-I, appendix-a103e7846d21-phase-6-\/-\/-source-analysis-and-verification)} · \hyperref[appendix-4ad47bdd269a]{Dependency Map}
